\subsection{The gateway holding the session}
\label{sec:ha:disabled:gateway}

The gateway carrying a member's session dies and is restarted. The member's
connection ends and it reconnects.

\subsubsection{What state it holds}
\label{sec:ha:disabled:gateway:state}

The gateway holds the conversation with the member: where the session's
numbering has reached in each direction, and which of the outbound numbers
carried execution reports rather than session-level traffic. That is the whole
of what matters when it dies.

While it is running it also holds the member's identity and the terms it was
admitted under, an authentication exchange part-completed, a resend being
served, and reports held back while the numbering is being established. None of
these needs to survive: the first is re-established at the next logon, and the
rest are abandoned and asked for again.

It holds no order state. The book belongs to the matching engine and the
ordering to the sequencer. Where the gateway must act on a member's behalf ---
cancelling what the member has resting when its connection ends --- it reports
the disconnection and the terms it was given to the component that owns the
orders, rather than keeping a view of them to act on itself. The terms
accordingly travel with the session, and are remembered wherever the session's
other state is remembered.

\subsubsection{Where the copy that outlives it lives}
\label{sec:ha:disabled:gateway:durable}

The sequencer holds the numbering. Gateways report a session's position as it
advances and the sequencer remembers it, because the sequencer is the only
component that every gateway instance talks to. Execution reports are held in
the sequencer's write-ahead log, stamped with the session they belong to.

An order that has not been committed has no copy anywhere in the venue. Between
the member sending it and the sequencer committing it the only record inside the
venue is the gateway's, and when the gateway is what died, nothing can recover
it.

That is deliberate rather than a shortcoming. A durable record kept by the
gateway would be a write on the hot path, in front of the write-ahead log that
already exists, doubling the cost of every order to protect a window of
microseconds --- and it would be redundant, because a copy of the order exists
outside the venue in the hands of the member that sent it. The venue therefore
does not attempt to recover an uncommitted order. It guarantees instead that
such an order never takes effect afterwards, and that the member can establish
as much and send it again.

\begin{gap}{}
  The gateway keeps its own view of what each session has resting, and acts on
  that view when a connection ends. The view is built only as execution reports
  arrive, so a restarted gateway holds nothing for orders placed before the
  restart, and a member that asked for cancel-on-disconnect is no longer
  protected by it, is not told, and cannot tell.
\end{gap}

\begin{requirement}{R-0008}{Reports remain recoverable for the reconnection window}
  The venue shall retain every execution report for at least as long as a member
  may take to reconnect and ask for it, and that period shall be stated.
  \rationale{Recovery that depends on a record still being present is only as
  good as the retention behind it. A retention period that is a by-product of
  storage rather than a stated figure cannot be relied on by anyone.}
  \uncovered{retention is not measured}
\end{requirement}

\begin{gap}{BUG-0048, BUG-0046}
  There is no stated retention period. What survives is whatever the
  write-ahead log has not yet discarded, which varies with the rate orders
  arrive at, so the window narrows exactly when a member has most at stake.
\end{gap}

\subsubsection{How a restarted instance becomes current}
\label{sec:ha:disabled:gateway:recovery}

The member reconnects, the session binds to the restarted gateway, and the
sequencer supplies the numbering it remembers. What comes back with it depends
on how far each order had travelled.

There is only one such operation, and that is worth stating because it is not so
elsewhere. A gateway is never promoted: it elects nothing and holds no role. A
member returning to the instance that restarted and a member arriving at a
different instance are therefore the same event seen from the venue --- a
session binds to an instance that does not hold its state, and that instance
asks the sequencer for it. Nothing in the sequence is particular to running more
than one gateway, which is why gateways need no arbiter, and why running several
of them costs so little.

\textbf{An order the gateway had received but not forwarded} was never
committed.

\begin{requirement}{R-0001}{An order that was not forwarded never takes effect}
  The order shall have no effect on the venue. No execution report shall be
  produced for it, and no order shall rest on the book on its behalf.
  \rationale{The venue had not taken the order, so acting on it later would
  execute an instruction the member has been given no reason to believe is
  live.}
  \uncovered{no scenario restarts a gateway}
\end{requirement}

\textbf{An order the sequencer had committed} is the venue's responsibility from
that moment.

\begin{requirement}{R-0004}{A committed order is processed}
  The order shall be processed by the matching engine, whether or not the
  gateway that received it is running when the engine reaches it.
  \rationale{Commitment is the venue taking the order. Nothing downstream of it
  may be conditional on the component that happened to carry it.}
  \uncovered{scenario 23 covers a gateway that is not restarted}
\end{requirement}

\textbf{A report that had been given an outbound sequence number} but did not
reach the member leaves a gap in the member's numbering.

\begin{requirement}{R-0007}{A numbered report that was not received is recoverable}
  The member shall detect the gap in its inbound numbering, and the venue shall
  supply the reports that occupied the missing numbers.
  \covers{ha_test:22}
\end{requirement}

\subsubsection{What the member is told}
\label{sec:ha:disabled:gateway:member}

\begin{requirement}{R-0002}{A member can establish that an order was not taken}
  On reconnecting, a member shall be able to determine that the venue holds no
  order for a \texttt{ClOrdID} it had sent.
  \rationale{The member cannot distinguish this case from a lost reply. Every
  venue provides an order status enquiry for exactly this reason.}
  \uncovered{the venue does not answer order status enquiries}
\end{requirement}

\begin{gap}{}
  The venue answers no order status enquiry of any kind, so a member has no way
  to ask what is held for it. Recovery rests entirely on what the venue chooses
  to send.
\end{gap}

\begin{requirement}{R-0003}{An order that was not taken may be resubmitted unchanged}
  A member may send the order again under the same \texttt{ClOrdID}, and the
  venue shall process it as a new order.
  \rationale{Without this the member must choose between abandoning an order and
  risking a duplicate. Resubmission under the original identifier is what makes
  the choice unnecessary.}
  \uncovered{no scenario restarts a gateway}
\end{requirement}

\begin{requirement}{R-0006}{A resubmitted committed order is refused as a duplicate}
  Where a member sends again an order the venue had already committed, under the
  same \texttt{ClOrdID}, the venue shall refuse it as a duplicate and shall not
  create a second order.
  \rationale{This is what makes \reqref{R-0003} safe. A member that cannot tell
  the two cases apart relies on the venue telling them apart on its behalf.}
  \uncovered{no scenario restarts a gateway}
\end{requirement}

\begin{requirement}{R-0010}{A duplicate is refused whatever has become of the first order}
  The refusal required by \reqref{R-0006} shall not depend on the order having
  reached the matching engine, nor on the engine that received it still running.
  \rationale{Commitment is the sequencer's act, so a record of what has been
  committed is the only thing that can answer this at the moment the question is
  asked. An order committed and not yet applied is invisible to the engine, and
  an engine that has restarted knows of no earlier submission at all; in either
  case a member following \reqref{R-0003} would be given a second live order.}
  \uncovered{no scenario restarts a gateway}
\end{requirement}

\begin{requirement}{R-0005}{A report produced for an unbound session is delivered on reconnection}
  An execution report produced while the session had no connection shall be
  delivered when that session next binds to a gateway, without the member having
  to ask for it.
  \rationale{No outbound sequence number was allocated for such a report, so it
  leaves no gap. The member has nothing to notice and nothing to request, and a
  resend cannot reach it.}
  \uncovered{no scenario restarts a gateway}
\end{requirement}

\begin{gap}{}
  Recovery is driven entirely by the member noticing a gap and asking. A report
  produced while the session had no connection was never numbered, so it leaves
  no gap, and nothing pushes it.
\end{gap}

\begin{requirement}{R-0009}{A gap that cannot be filled is reported as such}
  Where the venue cannot supply a report the member has asked for, it shall say
  so rather than fill the number silently.
  \rationale{A silently filled gap leaves the member believing it is
  synchronised when it has lost a report it will never see.}
  \uncovered{no scenario restarts a gateway}
\end{requirement}
